\section{Introduction}

Understanding surgical videos is fundamental to workflow recognition, context-aware assistance, retrospective review, quality assessment, and surgical education. Yet the labels required by supervised systems, including phases, steps, actions, tools, and triplets, remain expensive to annotate at scale. This bottleneck has made surgical video-language pretraining (VLP) increasingly attractive: instead of learning each downstream task from scratch, a model can first learn transferable surgical semantics from weakly paired video-text data and then support zero-shot or low-supervision evaluation.

Recent surgical VLP systems have shown that this strategy works. Lecture videos, narrated procedures, and curated corpora such as SurgLaVi provide scalable supervision and have enabled steady progress in surgical representation learning~\cite{lecturevideo2023,peskavlp2024,surglavi2025}. However, most existing pipelines still inherit a standard CLIP-style assumption: once a text span is paired with a video interval, the sampled clip is treated as a reliable positive visual unit for the paired text. Multiple frames are sampled from the annotated interval, pooled into one visual representation, and aligned to the text through a symmetric contrastive loss.

This assumption is convenient, but in surgical data it is often inaccurate. Our text supervision is typically not a manually curated caption for a cleanly trimmed clip. Instead, it is derived from narrations, transcripts, or weak temporal descriptions attached to a broader interval. In such settings, the textual evidence can be temporally offset from the most informative visual moment: a surgeon may describe the next action before it happens, summarize a broader sub-procedure after the key manipulation has passed, or narrate a local event whose discriminative visual evidence occupies only a short subsegment of the annotated interval. Similar weak synchronization has been reported in instructional and narrated videos beyond surgery~\cite{shen2021dwsa,jin2023ground_articles,xu2023procedure_aware_repr}. The result is a practical mismatch: a video-text pair may be globally valid while still containing only a small subset of frames that truly support the text.

This mismatch matters precisely because modern surgical benchmarks are not purely coarse-grained. Phase recognition can benefit from broad contextual cues, but step, action, tool, and triplet prediction rely much more heavily on localized visual evidence. If weak supervision is always consumed through uniform interval-level pooling, the representation may drift toward broad procedural context at the expense of the local cues required for fine-grained transfer. In other words, the problem is not simply noisy annotation; it is the training assumption that all frames inside a weakly paired interval are equally useful positives.

We argue that this problem should be addressed directly as \emph{temporal misalignment inside weak video-text pairs}. This view is consistent with ideas from multiple-instance learning and weak temporal localization, where a coarse positive unit is better treated as a partially positive bag than as a uniformly valid set of instances~\cite{dietterich1997mil,ilse2018attentionmil,li2023boostingwtal,hu2024mmplateau,li2024fps}. It is also aligned with recent video-language work showing that fixed temporal pooling often mixes relevant and irrelevant evidence, and that text-guided aggregation can recover more informative local correspondence~\cite{gorti2022xpool,li2023prost,liu2023stan}. In surgery, this observation becomes particularly important because temporal transitions are subtle, workflow descriptions are semantically dense, and narration-style supervision is rarely perfectly synchronized with visual events.

Based on this perspective, we present \textbf{SurgAlign}, a misalignment-aware training framework for surgical VLP. SurgAlign treats each weak interval as a noisy temporal anchor rather than a trustworthy clip boundary. During training, it expands the interval into a bounded temporal neighborhood, samples candidate frames from the expanded window, and uses the paired text to softly reweight local visual evidence. The resulting selected view is trained jointly with a standard global branch, so the model retains a stable global alignment signal while also learning to recover text-relevant local evidence from weakly aligned intervals. Crucially, we use multi-level annotations to condition this evidence-selection process itself. Fine, mid, and coarse captions differ not only in temporal scope but also in internal semantic complexity: fine captions are often dominated by a single cue, whereas broader captions more often behave as compositional semantic structures. SurgAlign therefore uses hierarchical labels to jointly control how many latent semantic cues may be extracted from the text and how sharply the temporal evidence is concentrated. In the hierarchical setting, we additionally regularize selected views from same-video adjacent levels so that local evidence remains compatible across nested temporal scopes.

This design is deliberately lightweight. SurgAlign does not introduce a new test-time retrieval protocol, a new surgical foundation model, or a heavy query-conditioned matching module. All additional logic is used only during pretraining, and inference remains standard CLIP-style independent video and text encoding. This makes the method easy to position: our contribution is not a new backbone, but a better way to consume weak temporal supervision on top of existing backbones.

The paper is organized around two complementary empirical questions. First, can misalignment-aware temporal denoising improve surgical VLP when compared under a strong surgical-domain backbone? Second, does the same training rule remain useful when the initialization backbone changes? We therefore use a surgical-domain primary model for the main benchmark comparison, and separately provide a cross-backbone validation to show that the gain is not tied to a single initialization family.

In summary, this paper makes three contributions. First, it identifies temporal misalignment inside weakly paired surgical video-text samples as an under-modeled problem in current surgical VLP. Second, it proposes SurgAlign, a training framework that combines expanded temporal evidence with multi-level annotation-aware evidence selection, where hierarchy jointly conditions latent semantic cue extraction and temporal concentration while preserving standard CLIP-style inference. Third, it validates the method both in a primary surgical-domain setting and in a cross-backbone setting, supporting the claim that the contribution lies in the training strategy rather than in a particular backbone choice.
